\question\TFQuestion{F}{Complex organisms always have more chromosomes than simple organisms. }
\question\TFQuestion{T}{Complex organisms can have fewer chromosomes than simple organisms. }

\question\TFQuestion{F}{Thomas Robert Malthus wrote \emph{Systema Natur\ae} that influenced Darwin.}
\question\TFQuestion{T}{Thomas Robert Malthus wrote the \emph{Principle of Populations} that influenced Darwin.}

\question\TFQuestion{T}{Alfred Russel Wallace also developed a concept of natural selection.}

\question\TFQuestion{F}{Darwin used his principles of natural selection to explain how artificial selection works in show animals, farm animals, and crops.}

%% Microevolution
\question\TFQuestion{F}{The smallest unit of evolution is the individual.}

\question\TFQuestion{T}{The smallest unit of evolution is the population.}

\question\TFQuestion{F}{Gene flow is enough to stop the effects of genetic drift in a small population.}
\question\TFQuestion{F}{Gene flow is enough to stop the effects of genetic drift in a large population.}

\question\TFQuestion{F}{Mutation is enough to stop the effects of genetic drift in a small population.}

\question\TFQuestion{T}{Microevolution and macroevolution are caused by the same underlying process: change in genetic variation in a population over time.}

% Genetic Drift

\question\TFQuestion{F}{Genetic drift is strongest in large populations.}

\question\TFQuestion{T}{Genetic drift is weakest in large populations.}

\question\TFQuestion{T}{Genetic drift is strongest in small populations.}

\question\TFQuestion{F}{Genetic drift is weakest in small populations.}

% Analogy vs homogy
\question\TFQuestion{F}{Homologies evolve by natural selection. Analogies do not.}
\question\TFQuestion{F}{Analogies evolve by natural selection. Homologies do not.}
\question\TFQuestion{T}{Both homologies and analogies can evolve by natural selection.}
\question\TFQuestion{F}{Analogies can be used to test hypotheses if two organisms do \emph{not} share a common ancestor.}


%% Natural Selection

\question\TFQuestion{F}{Adaptations are the result of descent with modification.}

\question\TFQuestion{T}{Adaptations can lead to descent with modification.}

\question\TFQuestion{F}{Natural selection acts on the population.}

\question\TFQuestion{T}{Natural selection acts on the individual.}

\question\TFQuestion{T}{If a trait is energetically costly and provides no survival advantage, then natural selection may cause that trait to disappear from the population over many generations.}

\question\TFQuestion{F}{If a trait is energetically costly and provides no survival advantage, then genetic drift will cause that trait to disappear from the population over many generations.}

\question\TFQuestion{T}{Genetic variation among individuals in a population is necessary for natural selection to act upon individuals in that population.}

\question\TFQuestion{F}{Overproduction of offspring by every individual in a population is necessary for evolution to occur.}

\question\TFQuestion{F}{The effects of mutations on the evolutionary success of organisms are always unfavorable.}

\question\TFQuestion{F}{Natural selection by genetic drift is most common in small populations.}

\question\TFQuestion{F}{Natural selection by genetic drift is most common in large populations.}

\question\TFQuestion{F}{Natural selection is a random process.}


%% Hardy Weinberg

\question\TFQuestion{F}{One assumption of the Hardy-Weinberg equations is that all individuals have an equally likely chance of migrating to another population.}

\question\TFQuestion{T}{One assumption of the Hardy-Weinberg equations is that individuals are not migrating into or out of the population.}

\question\TFQuestion{F}{One assumption of the Hardy-Weinberg equations is that mutations are random.}

\question\TFQuestion{T}{One assumption of the Hardy-Weinberg equations is that mutations do not occur.}

\question\TFQuestion{T}{One assumption of the Hardy-Weinberg equations is that natural selection is not present in a population.}


%% heterozygosity, homozygosity.

\question\TFQuestion{T}{Genetic drift always reduces heterozygosity in a population.}

\question\TFQuestion{F}{Genetic drift always reduces homozygosity in a population.}

\question\TFQuestion{F}{Genetic drift always increases heterozygosity in a population.}

\question\TFQuestion{T}{Genetic drift always increases homozygosity in a population.}

\question\TFQuestion{F}{Genetic drift does not change heterozygosity in a population.}

\question\TFQuestion{F}{Genetic drift does not change homozygosity in a population.}


\question\TFQuestion{T}{Non-random mating always reduces heterozygosity in a population.}

\question\TFQuestion{F}{Non-random mating always reduces homozygosity in a population.}

\question\TFQuestion{F}{Non-random mating always increases heterozygosity in a population.}

\question\TFQuestion{T}{Non-random mating always increases homozygosity in a population.}

\question\TFQuestion{F}{Non-random mating does not change heterozygosity in a population.}

\question\TFQuestion{F}{Non-random mating does not change homozygosity in a population.}



%% For the following T/F, consider giving p2 and asking about total homozygosity / heterozygosity. Also, make versions that ask about heterozygosity.

\question\TFQuestion{F}{If $p=0.4$, then total homozygosity in the population is 0.16 (assume equilibrium).}

\question\TFQuestion{T}{If $p=0.4$, then total homozygosity in the population is 0.52 (assume equilibrium).}

\question\TFQuestion{F}{If $p=0.4$, then total homozygosity in the population is 0.4 (assume equilibrium).}

\question\TFQuestion{T}{If $p=0.4$, then total heterozygosity in the population is 0.48 (assume equilibrium).}

\question\TFQuestion{F}{If $p=0.4$, then total heterozygosity in the population is 0.32 (assume equilibrium).}

\question\TFQuestion{F}{If $p^2 = 0.16$, then total homozygosity in the population is 0.16 (assume equilibrium).}

\question\TFQuestion{F}{If $p^2 = 0.16$, then total homozygosity in the population is 0.84 (assume equilibrium).}

\question\TFQuestion{F}{If $p^2 = 0.16$, then total homozygosity in the population is 0.32 (assume equilibrium).}

\question\TFQuestion{T}{If $p^2 = 0.16$, then total homozygosity in the population is 0.52 (assume equilibrium).}


%% Sexual Selection
\question\TFQuestion{F}{One difference between intrasexual and intersexual selection is that intrasexual selection does not cause sexual dimorphism.}

\question\TFQuestion{F}{Both intrasexual and intersexual selection cause the evolution of sexual dimorphism.}

\question\TFQuestion{T}{Sexual dimorphism is most often the result of directional selection.}

\question\TFQuestion{F}{Sexual dimorphism is most often the result of stabilizing selection.}

\question\TFQuestion{F}{Sexual dimorphism is most often the result of disruptive selection.}

\question\TFQuestion{T}{Sexual selection usually causes a pattern of directional selection.}

\question\TFQuestion{T}{Intrasexual selection acts directly on relative fitness.}

\question\TFQuestion{F}{Intrasexual selection acts indirectly on relative fitness.}

\question\TFQuestion{T}{Intersexual selection acts directly on relative fitness.}

\question\TFQuestion{F}{Intersexual selection acts indirectly on relative fitness.}

\question\TFQuestion{T}{Sexual selection usually causes a directional mode of selection.}


\question\TFQuestion{F}{Theories are tentative explanations that explain specific observations.}

\question\TFQuestion{F}{Microevolution is genetic change in individuals over time}. 
\question\TFQuestion{T}{Evolution is genetic change in a population over time.}


%% END

%% Consider adding mutation types (frameshift, point, etc.) to the relevant study guide.
\question\TFQuestion{T}{A mutation that changes a G to a T in \textsc{dna} is called a point mutation.}